\section{Operational and theorem-level comparisons}\label{sec:comparison}
The comparisons below fix the minimized object before comparing results.
They do not use a change of terminology as evidence of originality.

\subsection{Ordered randomized programs}
\begin{proposition}[The layered-program interpretation]\label{prop:program}
For a fixed command alphabet, Definition~\ref{def:machine} is equivalent
to a nonuniform probabilistic ordered read-once layered program with
cut-dependent widths and real terminal columns. For one selected column,
randomized Boolean acceptance has probability $(1+m)/2$, where $m$ is
the column's mean, without changing any command-layer width. A separate
charged terminal layer needs at most two labels. Maximum width is not
total program size, and neither measure prices the row descriptions.
\end{proposition}
\begin{proof}
At layer $t$, the outgoing edge weights labelled by command $a$ are
exactly the row $T_{t,a}(s,\cdot)$. Summing path probabilities gives
\eqref{eq:machine}. Conversely the edge weights of a probabilistic
layered program form those stochastic matrices. At the terminal label
$s$, toss a Bernoulli coin of parameter $(1+d_j(s))/2$. This implements
Boolean acceptance with the claimed probability. Choosing a query selects
one column; it does not request independent outcomes for all columns.
Padding a layer by unreachable labels and assigning normalized unused
rows proves the same assertion for varying widths.
\end{proof}
The equivalence permits arbitrary real edge probabilities and externally
known layers. The Fourier-growth results for regular and width-three
branching programs \cite{RSV,SVW,LPV} have different structural assumptions
and objectives. Here the external group average has a gap; the hidden
transition matrices need not be regular or have a gap. In particular,
a stochastic reset matrix is permitted. The all-spectra occupation
inequality is a lower bound for the specified matrix-coefficient response,
not a pseudorandom generator or a general lower bound for Boolean OBDDs.

\subsection{Zero-delay coding and state aggregation}
The structural coding results of Witsenhausen and Walrand--Varaiya,
as formulated in \cite[Theorems 1--2]{WLY}, reduce a Markov-source encoder
to the current source and a posterior based on past transmitted messages.
Wood--Linder--Y\"uksel prove stationary optimality and periodic finite-memory
approximation for average distortion \cite[Theorems 3--4]{WLY}. A posterior
is generally a point of a continuous simplex, not a finite alphabet.
Those results optimize ongoing reproduction for a given stochastic
source. Here a machine sees commands rather than the true orbit vector,
and every command word must satisfy a final numerical specification.
Only its current label, not a free transcript of sent messages, is retained.
Thus a sufficient belief state in the coding problem does not give a
cardinality upper bound in the present problem. Conversely our width
lower bound does not lower-bound the optimal average distortion of those
coders.

Nonanticipative rate-distortion minimizes directed information over
causal reproduction kernels under expected distortion constraints.
For example, the finite-horizon Markov-source analysis in
He--Charalambous--Stavrou \cite[Sections II--III]{HCS} derives a dynamic
optimization of those kernels and an alternating computation. That
objective is neither $\max_t K_t$ nor $\sum_tK_t$. In particular,
mutual information can be small for a channel whose exact finite positive
realization needs many labels. A source-law rate-distortion optimization
cannot be inserted in place of the wordwise specification without a
separate reduction.

Geiger--Temmel \cite[Theorem 9]{GT} characterize higher-order strong
lumpability by an equality of conditional entropies for a stationary
finite Markov chain and a deterministic lumping. Their entropy-rate
preservation results concern the whole lumped process. We allow
randomized, epoch-dependent maps, retain a command-driven terminal
observable, and do not require preservation of the trajectory entropy
rate or of a Markov partition. State aggregation is therefore an adjacent
comparison, not a theorem claimed to have been subsumed. The common
use of conditional expectations does not identify the optimization tasks.

\subsection{Conditional expectations and Gaussian entropy}
Atalik--K\"ose--Gastpar \cite[Theorem 1 and Proposition 2]{AKG} study
$M=\E[X\mid X+W]$ for Gaussian observation noise. In the identity
observation case their inequality has the form
$\Ent(M)\ge2\Ent(X)-\Ent(X+W)$, with a vector extension. Lemma
\ref{lem:centroid} instead permits an arbitrary finite conditioning
variable $J$, adds new independent noise \emph{after} conditioning, and
upper-bounds a difference of two smoothed entropies by the martingale
coupling's squared displacement. A discrete $M$ is allowed here, whereas
its unsmoothed differential entropy is not the quantity being compared.
The two displayed inequalities are not direct substitutions of one
another. The Gaussian score and posterior-covariance identities used in
\eqref{eq:hessian} are classical, and our short semiconcavity proof is
not claimed as a new estimation identity. Polyanskiy--Wu's Wasserstein
entropy continuity \cite{PW} provides a further comparison; the quadratic
cost here uses the particular conditional-mean coupling, not arbitrary
transport. The argument needed for the occupation theorem is the
amortization of that cost together with orbital sparsity.

\subsection{Positive realization and orbit geometry}
Positive realization and invariant convex sets have a long history
\cite{Heller,Vidyasagar}. The occupation theorem does not require an
internal state to be a point on an observed orbit or in the residual
span. Such a restriction would miss hidden stochastic lifts. Classical
sphere and projective covering exponents explain the net counts; the
additional issue for conjugation is uniformity over spectra of all
conditional centroids. Theorem~\ref{thm:all-spectra} settles that issue
through a separated spectral projector and the explicit Grassmann bound.
Its exponent is matched by an exact projective, rather than ambient
spherical, construction. This is the reason that $q-1$, not
$(q^2-2)/2$, appears in Theorem~\ref{thm:main}.

The deep group gaps are external inputs. Proposition~\ref{prop:algebraic-gates}
uses a qualitative theorem; Proposition~\ref{prop:LPS} uses a numerical
Hecke norm theorem for a different named alphabet. Finite diagonalization
of a few harmonics is not a proof of either full gap. Neither density of
a subgroup nor a gap in its defining finite-dimensional representation
is silently promoted to a full $L^2$ norm gap.

\subsection{Implementation and the remaining distinctions}
All exact theorems here count atomic nonuniform labels. Replacing each
row by a rational approximation and then using a fixed-duration padded
fair-bit comparison sampler, as in the retained predecessor, incurs
additional control labels and terminal approximation error. The padding
removes dependence of duration on private coin outcomes; it does not
price the read-only threshold tables or prove uniform space complexity.
No exact finite-coin compiler for arbitrary irrational rows is asserted.

The proved results are the representation-valued occupation inequality,
the all-spectra conjugation estimate, the matched projective cardinality
order, and the effective six-command consequence. They do not classify
all compact representations, all finite alphabets, or arbitrary adaptive
experimental design. The sufficient calibration threshold is not shown
optimal, and signal approaching zero with the horizon requires retaining
the explicit constants rather than invoking a fixed-signal asymptotic.
Exact finite integer optima and leading constants remain different
questions. The standalone article contains all proofs used for its new
claims; preservation records and computation are not mathematical or
priority certificates.
